\chapter{Verification Methodology, Compared Across Platforms}
\label{ch:verification-methodology}

Chapter~\ref{ch:testing-philosophy} in this Part covers the individual verification
\emph{mechanisms} this ecosystem relies on --- pixel tests, differential testing, the D3D9
oracle, sanitizers, the \texttt{NOXNA} build gate. This short closing chapter of Part~VIII takes
a different cut through the same material: comparing how much of that machinery actually
reaches each of the four platforms the previous three chapters covered, side by side, because
the honest answer varies more than a single ``cross-platform and verified'' claim would
suggest.

\section{A comparative summary}

\begin{center}
\small
\begin{tabular}{lp{8.7cm}}
\toprule
\textbf{Platform} & \textbf{What is actually verified, as of this book's source material} \\
\midrule
Linux (native) & The deepest verification of any platform: the full roughly~4{,}370-test unit
  suite, the roughly~490-test GPU pixel suite across four backends, real differential testing
  against a genuine FNA install (Chapter~\ref{ch:windows-wine}), and the sanitizer trio where
  applied (Chapter~\ref{ch:testing-philosophy}). This is the ground every other platform's own
  claims are ultimately measured against. \\
Windows (via Wine) & Real GPU-accelerated rendering through DXVK/vkd3d-proton, gated by a
  self-verifying marker check so a silent WineD3D fallback cannot corrupt a result
  (Chapter~\ref{ch:windows-wine}); includes the single strictest test in the whole project,
  the D3D9 oracle against a genuinely running XNA~4.0 install. Real Windows hardware
  verification remains open on all three Direct3D backends. \\
Web (Emscripten) & Split sharply by backend. \texttt{CANVAS}: a real build exists, links, and
  runs its test suite under Node --- but no test in this environment can touch a real
  \cnaclass{GraphicsDevice} at all, since Node provides no WebGL/Canvas context, leaving a
  13-item manual browser checklist entirely unchecked (Chapter~\ref{ch:web-emscripten}).
  \texttt{EASYGL}/WebGL2, the platform's actual default backend, has never been built under
  Emscripten at all in this project's history. \\
Android (NDK) & Split sharply by demo, not just by backend. The devices/sensors demo: real,
  live-verified on an x86\_64 emulator, down to synthetic sensor injection producing a real
  on-screen response. Graphics (\texttt{SDL\_RENDERER}, Android's own default): never run on
  any device or emulator, and, per a documented regression in an unrelated sibling project,
  currently not even confirmed to \emph{build} (Chapter~\ref{ch:android-ndk}). \\
macOS & The one platform in this comparison with no dedicated chapter of its own in this
  book, and worth naming as exactly that kind of gap rather than silently omitted: the
  project's own input-platform notes confirm mouse warp and global cursor position genuinely
  work through the SDL3 Cocoa backend, the same as X11, but explicitly flag every macOS
  specific as ``manual-gated,'' since the CI matrix is Linux-only and macOS is
  ``not headless-verifiable'' there at all --- no CI, no emulator equivalent, and no
  automated substitute the way Wine stands in for Windows. \\
\bottomrule
\end{tabular}
\end{center}

\section{Why the gradient exists, and why it is not a criticism}

None of the gaps this table names are hidden anywhere in this ecosystem's own documentation
--- every one of them was found, in this book's own research, stated plainly in the project's
own docs, exactly the way Chapter~\ref{ch:what-is-cna} described this project's
overall relationship with its own claims. The gradient itself is a predictable consequence of
where the actual development happens: this is a Linux development environment first, with
Windows reached through an elaborate but real Wine-based bridge, and Web/Android reached
through cross-compilation toolchains that can verify \emph{building} far more easily than
they can verify \emph{running correctly}, absent a real browser or a real device. A native
Linux pixel test can assert on an exact color value; an Emscripten build that cannot even
construct a graphics device has no equivalent assertion available to it at all, no matter how
carefully the surrounding C\texttt{++} is written.

This book's own coverage has exactly the same shape of gradient, worth naming rather than
leaving implicit: Windows, Web, and Android each earned a full chapter in this Part because
each has a substantial, documented, mechanically-interesting verification story to tell; macOS
has real, working input support confirmed in the project's own platform notes, but no
comparably deep verification narrative exists to expand into its own chapter, since --- unlike
Windows' Wine bridge or Android's emulator --- there is no equivalent automated or
semi-automated stand-in for a real Mac in this project's own dev loop at all.

\section{What this means for a reader of this book}

If you are choosing a target platform for a new CNA project, or evaluating CNA for a port,
this table is the single fact most likely to change your plan: build success and behavioral
correctness are two different claims, verified to different degrees on different platforms,
and this ecosystem's own documentation is careful --- and this book has tried to be equally
careful --- to say which one it means every time it makes one. Treat a Linux-verified claim as
the strongest evidence available in this project today; treat a Windows claim as strong but
not hardware-final; and treat a Web or Android graphics claim as, at best, ``compiles,'' until
someone actually runs it where you intend to ship it.
